\subsection{Raw toric surgery does not descend to cut isotopy classes}
\label{sec:A1-toric-descent}

There remains a distinction between replacing a cut by its unique
geodesic and asserting that raw cutting and capping is independent
of its representative. The latter is suggested by the discussion
of a fixed welded $X$ and nearby seams in 67/079--080. It is false
even if all the representatives under consideration are torically
compatible. We give an exact analytic counterexample; its phase
is specified by a deterministic analytic rule, not by an evaluated
numerical angle.

\begin{lemma}[Proved: the third jet of a one-active-cap pant]
\label{A1-pant-third-jet}
Let the hyperbolic pant have one boundary of length $2\pi\lambda$,
$\lambda>0$, and two cusps. Normalize its capped sphere by placing
the cap centre at $0$ and the ordered cusps at $1,\infty$.
Let $r(\lambda)$ be the cap's conformal radius, and normalize its
unit-disc coordinate $u_\lambda$ to have positive derivative at
zero. Then
\[
 g_\lambda(z):=r(\lambda)u_\lambda(z)
       =z+\frac12z^2+\beta(\lambda)z^3+O(z^4),
 \qquad
 \beta(\lambda)=\frac5{16}
                 +\frac1{16(\lambda^2+4)}.
\]
Moreover $r:(0,\infty)\to(0,4)$ is a strictly increasing
real-analytic bijection.
\end{lemma}
\begin{proof}
The radius assertion is the one-active-cap case of
Theorem~\ref{A2b-image} and the Koebe endpoint in
Corollary~\ref{A2-explicit-limits}. Put
$a=(1+\lambda^2)/2$. We use the projective-connection argument
of Proposition~\ref{A2-schwarzian} with its cusp extension made
explicit. At a cusp a developing coordinate is a nonzero constant
multiple of $\log q$, whose Schwarzian is $1/(2q^2)$.
Thus the double-pole coefficients at $1,\infty$ are both $1/2$.
At the positive boundary they are $(1+\lambda^2)/2$ as before.
The unique projective connection on the three-pointed sphere with
these principal coefficients is consequently
\[
 T(z)=\frac a{z^2}+\frac1{2(z-1)^2}
                         -\frac a{z(z-1)}
      =\frac a{z^2}+\frac a z+(a+1/2)+O(z).
\]
The cap developing map gives
$T=a(g'/g)^2+\{g,z\}$. For
$g=z+bz^2+\beta z^3+O(z^4)$, its coefficients of $z^{-1}$ and
$z^0$ are respectively
\[
 2ab,\qquad (4a+6)\beta-(a+6)b^2.
\]
Matching the displayed connection gives $b=1/2$ and
\[
 \beta=\frac{5a+8}{16a+24}
       =\frac5{16}+\frac1{16(\lambda^2+4)}.
\]
The normalization by $r$ changes neither logarithmic derivatives
nor the Schwarzian, so this calculation applies to the stated
normalized germ.
\end{proof}

\begin{lemma}[Proved: a four-point plumbing coefficient]
\label{A1-four-point-plumbing-jet}
On two spheres with node $0$ and external markings $1,\infty$,
take node-coordinate germs
\[
 g_i(z)=z+\frac12z^2+\beta_i z^3+O(z^4),\qquad i=1,2.
\]
Plumb by $g_1(z)g_2(w)=t$. On the resulting four-pointed sphere,
normalize the global coordinate so that the first component's
$\infty$ is $\infty$, its $1$ is $1$, and the second component's
$\infty$ is $0$. Let $\Lambda(t)$ be the coordinate of the second
component's $1$. Then
\[
 \Lambda(t)=t-\frac12t^2+
 \left(\frac9{16}-\frac54(\beta_1+\beta_2)
                         +4\beta_1\beta_2\right)t^3+O(t^4).
\]
In particular, for the germs of Lemma~\ref{A1-pant-third-jet},
the cubic coefficient is
\[
 C(\lambda_1,\lambda_2)
   =\frac{11}{64}
       +\frac1{64(\lambda_1^2+4)(\lambda_2^2+4)}.
\]
\end{lemma}
\begin{proof}
For this coefficient calculation choose sufficiently small fixed
plumbing discs whose closures lie in $|z|<1/2$ and $|w|<1/2$.
This is possible for any of the coordinate germs in the statement.
The retained pieces then contain the circles $|z|=1$, $|w|=1$
and both external markings. First normalize a meromorphic
coordinate $Z$ to have its sole simple
pole at the first $\infty$, leading term $z$ there, and value zero
at the second $\infty$. On the two retained pieces write
\[
 Z_1(z,t)=z+B(t)+\sum_{n\geq1}C_n(t)z^{-n},\qquad
 Z_2(w,t)=\sum_{n\geq1}D_n(t)w^{-n}.
\]
These are convergent Laurent expansions for sufficiently small
$t$, in particular at $z=w=1$; their coefficients are meromorphic
on the whole central components. They depend holomorphically
on $t$, including $t=0$, where
$Z_1=z$ and $Z_2=0$. One can see this directly from the family of
meromorphic functions with the specified simple pole: on the
central genus-zero nodal curve its pole divisor has multidegree
$(1,0)$, with vanishing first cohomology, and evaluation and the
leading coefficient give the two normalizations. Equivalently,
Laurent principal-part matching in the plumbing annulus constructs
the same functions.

In the local total space $xy=t$, holomorphicity at the node and
$Z_2(w,0)=0$ show that
\[
 C_n(t)=O(t^{n+1}),\qquad D_n(t)=O(t^n).
\]
Indeed a term $y^m$ on the second branch has coefficient vanishing
at least once in $t$; after $y=t/x$ it contributes at least
$t^{m+1}$ to a pole of order at most $m$ on the first branch.
The analogous first-branch terms give the stated bound for $D_n$.

For completeness all terms needed for the cubic coefficient are
as follows. Set $A_i=\beta_i-1/4$ and $a=1/2-\beta_1$.
The inverse germ gives, on the overlap,
\[
 z=\frac t{g_2(w)}-\frac12\frac{t^2}{g_2(w)^2}
                          +a\frac{t^3}{g_2(w)^3}+O(t^4).
\]
Matching negative and constant Laurent coefficients at $w=0$
gives
\begin{align*}
 C_1(t)&=A_2t^2+O(t^3),\\
 B(t)&=\tfrac12t+(\tfrac12-\tfrac32\beta_2)t^2+O(t^3),\\
 D_1(t)&=t+\tfrac12t^2+
       \left[a(\tfrac32-3\beta_2)+A_1A_2\right]t^3+O(t^4),\\
 D_2(t)&=-\tfrac12t^2-\tfrac32a\,t^3+O(t^4),\\
 D_3(t)&=a\,t^3+O(t^4).
\end{align*}
For example, the coefficient of $w$ in $1/g_2(w)$ is
$1/4-\beta_2=-A_2$, which gives the first line. The negative
Laurent parts of $g_2^{-2}$ and $g_2^{-3}$ give the next three
principal-part lines; the $A_1A_2$ term comes from the coefficient
of $t/g_2(w)$ in $C_1(t)/z$. Higher $C_n$ cannot contribute to
these negative parts before order $t^4$.

It follows that
\begin{align*}
 Z_1(1,t)&=1+\tfrac12t+(\tfrac14-\tfrac12\beta_2)t^2+O(t^3),\\
 Z_2(1,t)&=t+\left[a(1-3\beta_2)+A_1A_2\right]t^3+O(t^4).
\end{align*}
The required cross-ratio is $\Lambda=Z_2(1,t)/Z_1(1,t)$.
Division gives the stated symmetric cubic coefficient. Finally
\[
 \frac9{16}-\frac54(\beta_1+\beta_2)+4\beta_1\beta_2
 =\frac{11}{64}
      +4(\beta_1-5/16)(\beta_2-5/16),
\]
which gives the last formula. This is a coefficient identity
derived from convergent plumbing, not a numerical approximation.
Enlarging the coordinate discs gives the same family wherever
both constructions are defined, by
Proposition~\ref{A1-welding-plumbing}; thus later larger parameter
domains give its holomorphic continuation and do not change
these Taylor coefficients.
\end{proof}

\begin{theorem}[Refuted: raw toric surgery is invariant under isotopy in a fixed surface]
\label{A1-toric-nondescent}
There is a four-labelled punctured sphere $X$ with two isotopic
essential analytic cuts $C,C_{\mathrm{geo}}$ such that both cuts
are canonically torically compatible, $C_{\mathrm{geo}}$ is the
unique hyperbolic geodesic in their isotopy class, and raw canonical
cutting and capping produces different nodal tensors. Thus raw
toric surgery does not descend to a functor on cut isotopy classes.
Replacing a representative by its geodesic before applying surgery
defines the different, well-defined functor already studied.
\end{theorem}
\begin{proof}
We specify all the choices. Use the radius function in
Lemma~\ref{A1-pant-third-jet}, and set
\[
 r_1=r(1),\qquad r_*=r(1/2),\qquad
 r_2=\frac{r_*^2}{4r_1},\qquad
 r_0=\frac{r_*}{2},\qquad R=r_0^2=r_1r_2.
\]
These are positive and less than $4$. Let
$\lambda_2=r^{-1}(r_2)$ and $\lambda_0=r^{-1}(r_0)$; strict
monotonicity gives $0<\lambda_2<\lambda_0<1/2$.
The reference pants defining $r_1,r_*$ have the explicit boundary
lengths $2\pi,\pi$. All other pants are uniquely fixed by the
displayed conformal radii, so no unchosen cap realization remains.

For the fixed pair $(\lambda_1,\lambda_2)=(1,\lambda_2)$ let
$\Lambda_U(t)$ be the cross-ratio of its plumbing family from the
preceding lemma. For the fixed equal pair $(\lambda_0,\lambda_0)$
write $\Lambda_G(t)$. These are honest holomorphic families on a
disc containing $|t|=R$, not just formal germs. To make this
domain explicit, put
\[
 \delta(\lambda)=\arctan\!\bigl(\operatorname{csch}(\pi\lambda)\bigr),
 \qquad
 a(\lambda)=r(\lambda)
                 \exp\!\left(\frac{\delta(\lambda)}{2\lambda}\right).
\]
The collar extension used in Theorem~\ref{A2b-image} makes
$g_\lambda$ a valid node coordinate on an embedded disc of radius
$a(\lambda)>r(\lambda)$, avoiding both external markings.
The factor $1/2$ keeps this disc strictly inside the available
collar. Thus the usual plumbing quotients exist for
\[
 |t|<\min\{a(1)a(\lambda_2),a(\lambda_0)^2\}=R\rho,
 \qquad \rho>1.
\]
At $t=R\eta$, $|\eta|=1$, their cut along the circles
$|g_i|=r_i$ is precisely the actual toric welding
$u_1u_2=\eta$, by Proposition~\ref{A1-welding-plumbing}.

Define the holomorphic function on $|c|<\rho$
\[
                 \Delta(c)=\Lambda_U(Rc)-\Lambda_G(Rc).
\]
Its cubic coefficient is nonzero. Indeed, with
$h(\lambda)=1/(\lambda^2+4)$,
\[
 h(1)h(\lambda_2)<\frac1{20}
       <\frac{16}{289}<h(\lambda_0)^2,
\]
and Lemma~\ref{A1-four-point-plumbing-jet} gives
\[
 [c^3]\Delta(c)
   =\frac{R^3}{64}
       \left[h(1)h(\lambda_2)-h(\lambda_0)^2\right]
   <-\frac{31R^3}{369920}<0.
\]
Consequently $\Delta$ is not the zero function. It has only
finitely many zeros on the unit circle, since that circle is a
compact subset of its open holomorphy disc.

For a deterministic phase, define
\[
 N=\min\{n\geq1:\ \Delta(e^{i/n})\ne0\},
 \qquad \eta_*=e^{i/N}.
\]
This integer exists: otherwise the distinct zeros $e^{i/n}$ would
accumulate at $1$, an interior point of the holomorphy disc, and
the identity theorem would contradict the nonzero cubic coefficient.
This is an exact analytic specification of the phase; no numerical
value of $N$ is asserted or needed. More generally every unit
phase outside the finite zero set gives the same counterexample
construction.

Now take $X$ to be the unequal-pair welding
$P(r_1)\cup_{u_1u_2=\eta_*}P(r_2)$, with its actual analytic
seam $C$ and its four distinct external labels. Its normalized
nodal output is the fixed pair of three-pointed spheres, with
tensor $q=R\eta_*$. The equal-pair welding
$X_G=P(r_0)\cup_{u_1u_2=\eta_*}P(r_0)$ has the same tensor,
and its seam is geodesic. But
$\Delta(\eta_*)\ne0$ says that $X$ and $X_G$ are not isomorphic
as four-labelled punctured spheres: their specified cross-ratios
differ.

Let $C_{\mathrm{geo}}$ be the unique hyperbolic geodesic of $X$
in the essential isotopy class of $C$. Raw cutting along this
geodesic is torically compatible. If its tensor were also
$R\eta_*$, the unique geodesic inverse in
Theorem~\ref{A1-graph-range}, or its explicit four-point case
Corollary~\ref{A1-small-ranges}, would identify its welded surface
with $X_G$. This contradicts the cross-ratio inequality. Thus
the two isotopic cuts in this same $X$ have different tensors,
which proves non-descent. No inference from the dimension of
the fixed-target fibers is used.
\end{proof}
